\begin{proof}[Proof of Theorem 55]
We construct an oracle $\mathcal O\in\mathsf{EXP}$ under which a dense subset of $\mathsf{MKtP}$ is in $\mathsf P$. We start with $n:=1$, $\mathcal O:=\emptyset$, and $F:=\emptyset$, where $F$ is a set of ``frozen'' strings. The construction of $\mathcal O$ in Stage $n$ is as follows. Let $\mathcal O_n$ denote the state of the oracle $\mathcal O$ at the beginning of stage $n$. Let $E_n$ be the set of all strings $x\in\{0,1\}^n$ such that $\mathrm{Kt}^{\mathcal O_n}(x)\leq n-c\log n$, where $c$ is a sufficiently large constant (e.g., $c:=3$). Let $F_n$ be the set of all strings $q\in\{0,1\}^*$ such that there exist $k\in[n-c\log n]$ and a description $d\in\{0,1\}^k$ of a Turing machine such that the universal Turing machine $U^{\mathcal O_n}$ on input $d$ makes the query $q$ in time $2^{n-k-c\log n}$. Then, we update $F:=F\cup F_n$ and $\mathcal O_{n+1}:=\mathcal O_n\cup(\{0,1\}^n\setminus(E_n\cup F))$ and move on to the next stage $n+1$. This completes the description of the oracle $\mathcal O:=\bigcup_{n\in\mathbb N}\mathcal O_n$. It is easy to observe that $\mathcal O\in\mathsf{EXP}$.

The oracle $\mathcal O$ is a dense subset of $\{x:\mathrm{Kt}^{\mathcal O}(x)>n-c\log n\}$ in the following sense:
\begin{enumerate}
\item $x\notin\mathcal O$ for every string $x\in\{0,1\}^n$ such that $\mathrm{Kt}^{\mathcal O}(x)\leq n-c\log n$.
\item $x\in\mathcal O$ for at least half of the strings $x\in\{0,1\}^n$.
\end{enumerate}
To see the first property, it suffices to show that if $x\in\{0,1\}^n$ satisfies $\mathrm{Kt}^{\mathcal O}(x)\leq n-c\log n$, then $\mathrm{Kt}^{\mathcal O_n}(x)\leq n-c\log n$ (which implies $x\in E_n$ and thus $x\notin\mathcal O$). Assuming that there exists a string $d$ such that $U^{\mathcal O}(d)$ outputs $x$ in time $2^{n-c\log n-|d|}$, we claim that $U^{\mathcal O_n}(d)$ also outputs $x$. If not, there exists a query in $\mathcal O\setminus\mathcal O_n$ made by $U^{\mathcal O}(d)$. Let $q$ be the first query in $\mathcal O\setminus\mathcal O_n$. The same query is also queried during the computation of $U^{\mathcal O_n}(d)$. Thus, $q\in F_n$, which implies $q\notin\mathcal O$, which is a contradiction.

To see the second property, we bound the size of $E_n\cup F$ at Stage $n$. The size of $E_n$ is at most $\sum_{k\in[n-c\log n]}2^k\leq2^n/n^{c-1}$. The number $|F_n|$ of frozen strings in Stage $n$ is at most
\[\sum_{k\in[n-c\log n]}2^k\cdot2^{n-k-c\log n}\leq2^n/n^{c-1}.\]
Thus, the number $|F|$ of frozen strings until Stage $n$ is at most $\sum_{n'=1}^n|F_{n'}|\leq2^n/n^{c-2}$. Overall, we obtain
\[|\mathcal O\cap\{0,1\}^n|\geq2^n-|E_n|-|F|\geq\tfrac12\cdot2^n.\]
Let $Q:=\mathcal O$. We prove that $Q\in\mathsf P^{\mathcal O}$ is a dense property such that every $2^n/n^{\omega(1)}$-time deterministic algorithm $A$ fails to find a string in $Q\cap\{0,1\}^n$ on input $1^n$. By the definition of Kt, we have $\mathrm{Kt}^{\mathcal O}(A(1^n))\leq O(\log n)+\log(2^n/n^{\omega(1)})\leq n-\omega(\log n)$. Thus, the output $A(1^n)$ of $A$ is not in $\mathcal O$.

We now use the worst-case to average-case reduction of [34] to obtain an approximation algorithm $A^{\mathcal O}$ for $\mathsf{Gap}_{\varepsilon}\mathsf{MKtP}$. Let $x\in\{0,1\}^n$ and $s\in\mathbb N$ be an instance of $\mathsf{Gap}_{\varepsilon}\mathsf{MKtP}$. Let $d=k\cdot O(\log^2 n)$, where $k=k(\varepsilon)$ is a sufficiently large constant that will be chosen depending on $\varepsilon>0$. We may assume without loss of generality that $O(k^2\cdot\log^2 n)\leq s\leq n+O(\log n)$ because whether $\mathrm{Kt}(x)\leq O(k^2\cdot\log^2 n)$ or not can be decided in time $n^{O(k^2\cdot\log n)}$ by an exhaustive search. Let $m$ be a parameter chosen later. Let $G$ be the pseudorandom generator construction of Lemma 56. The algorithm $A^{\mathcal O}$ accepts $(x,s)$ if and only if $G(x,z)\notin\mathcal O$ for every $z\in\{0,1\}^d$. The running time of $A^{\mathcal O}$ is $2^d\operatorname{poly}(n)=n^{O(\log n)}$.

We prove the correctness of $A^{\mathcal O}$. Assume that $\mathrm{Kt}^{\mathcal O}(x)\leq s$. Then, for every $z\in\{0,1\}^d$, we have
\[\mathrm{Kt}^{\mathcal O}(G(x,z))\leq\mathrm{Kt}^{\mathcal O}(x)+O(d+\log n)\leq s+O(d+\log n)\leq m-c\log m,\]
where the last inequality holds by choosing a sufficiently large $m=s+O(d+\log n)$. It follows that $G(x,z)\notin\mathcal O$ and that $A^{\mathcal O}$ accepts $(x,s)$. Conversely, assume that $A^{\mathcal O}$ accepts $(x,s)$, which means that $\Pr_{z\sim\{0,1\}^d}[\mathcal O(G(x,z))=1]=0$. We claim that $\mathrm{Kt}^{\mathcal O}(x)\leq(1+\varepsilon)\cdot s$. Since $\mathcal O$ is dense, we have
\[\Pr_{w\sim\{0,1\}^m}[\mathcal O(w)=1]-\Pr_{z\sim\{0,1\}^d}[\mathcal O(G(x,z))=1]\geq\tfrac12.\]
By Lemma 56, we obtain
\begin{multline*}
K^{\operatorname{poly}(n),\mathcal O}(x)\leq\exp(O(\log^2 n)/d)\cdot m+d+O(\log n)\\
\leq\exp(1/k)\cdot m+O(k\cdot\log^2 n).
\end{multline*}
In particular,
\begin{multline*}
\mathrm{Kt}^{\mathcal O}(x)\leq\exp(1/k)\cdot m+O(k\cdot\log^2 n)\\
\leq(1+O(1/k))\cdot m+O(k\cdot\log^2 n)\leq(1+\varepsilon)\cdot m,
\end{multline*}
where we choose a sufficiently large $k:=O(1/\varepsilon)$ and use that $m\geq O(k^2\cdot\log^2 n)$ in the last inequality.
\end{proof}
